\section{Conclusion}
\label{sec:conclusion}

We built and evaluated an isolated sign recognition system based on MediaPipe landmarks, under a
signer-independent protocol and with a fully client-side deployment. The three research questions
receive the following answers.

\paragraph{RQ1 Generalisation.} A 2.2\,M-parameter convolution--Transformer recognises
\numAslClasses{} ASL signs produced by unseen signers with \numTrTestTop{}\,\% top-1 accuracy, and
\numLsfbClasses{} LSFB glosses with \numLsfbScratchTest{}\,\%. A random split inflates the ASL
accuracy by \numLeakGap{} points, and an early-stopping artefact had overstated the effect of one
ablation fivefold: evaluation choices influence conclusions more than most design choices.

\paragraph{RQ2 Transfer.} Pre-training on ASL does not improve LSFB recognition at any
training-set size, even with five examples per sign, and with all training data it lowers test
accuracy by a small but significant margin. Differences in recording and elicitation
conditions likely outweigh what the two languages share at the level of pose features.

\paragraph{RQ3 Deployment.} The complete pipeline runs in a web browser, classifies a sign in
\numFpBrowser{}\,ms on one CPU thread and reproduces the predictions of the PyTorch model on every
tested clip. INT8 quantisation reduces model size but not latency on CPU.

Future work includes self-supervised pre-training on in-domain continuous signing, a rejection
class for out-of-vocabulary signs, richer non-manual features, and an evaluation of the demo with
Deaf signers.
